\documentclass[12pt,a4paper]{article}

\usepackage[utf8]{inputenc}
\usepackage[T1]{fontenc}
\usepackage{lmodern}
\usepackage{amsmath,amssymb,amsthm}
\usepackage[top=2.5cm,bottom=2.5cm,left=2.8cm,right=2.8cm]{geometry}
\usepackage{parskip}
\usepackage{booktabs}
\usepackage{xcolor}
\usepackage{colortbl}
\usepackage{hyperref}
\hypersetup{colorlinks=true, linkcolor=blue!60!black}

\newtheorem{theorem}{Théorème}[section]
\newtheorem{lemma}[theorem]{Lemme}
\theoremstyle{remark}
\newtheorem{note}{Note}[section]

\newcommand{\Inc}{\ensuremath{\mathrm{Include}}}
\newcommand{\Exc}{\ensuremath{\mathrm{Exclude}}}
\newcommand{\vbar}{\bar{v}}

\begin{document}

\begin{center}
  {\LARGE\bfseries Rôle de $S$ et $\alpha$ : théorie et validation (XOR, Iris)}
\end{center}
\vspace{0.5em}
\hrule
\vspace{1.5em}

\section{Notations}

$x\in\{0,1\}$ : feature ; $y\in\{0,1\}$ : label. $v,\vbar$ : automates des
littéraux $x$ et $\neg x$. $S,\alpha\in(0,1]$ : paramètres de la règle à 4 cas
(Section~2 de \textit{Démonstration de convergence}). Statistiques du canal de
bruit : $c=P(x{=}1)$, $a=P(y{=}1\mid x{=}1)$, $b=P(y{=}1\mid x{=}0)$.
$\rho=p^+/p^-$ : $\rho>1\Rightarrow$ l'automate converge vers \Inc ;
$\rho<1\Rightarrow$ vers \Exc{} (Lemme 6.4 du même document).

\section{Formule du seuil critique}

Pour le littéral pertinent (cas IDENTITY, $a>b$), le document précédent établit :
\[
  p^+_v = (1-c)(1-b)\,S, \qquad p^-_v = (1-c)\,b\,\alpha + c(1-a)\,S.
\]
$\rho_v>1 \iff p^+_v > p^-_v$. En résolvant pour $\alpha$ à $S$ fixé :
\[
  \boxed{\alpha^\star(S) = \frac{\Delta\,S}{(1-c)\,b}, \qquad
  \Delta = (1-c)(1-b) - c(1-a).}
\]
Pour $\alpha<\alpha^\star(S)$ : $\rho_v>1$, le littéral apprend correctement.
Pour $\alpha>\alpha^\star(S)$ : $\rho_v<1$, il apprend dans le mauvais sens.

\section{Mesure de $c$, $a$, $b$ sur XOR}

Sur \texttt{NoisyXORTrainingData.txt} (12 features, 5000 exemples) :
\begin{enumerate}
  \item Chaque feature seule : $P(x_i{=}1)\approx0{,}5$ et $P(x_i{=}y)\approx0{,}5$
  --- aucune ne prédit $y$ seule (attendu pour un XOR).
  \item Chaque paire $(i,j)$ : $P(x_i\oplus x_j = y)$. Une seule paire sort du
  lot, $(1,2)$, à $60{,}3\,\%$ --- toutes les autres $\approx50\,\%$.
  \item Donc bruit $=1-0{,}603=0{,}397$. En conditionnant sur $x_1$, $x_2$
  devient le littéral pertinent : $c=0{,}5$, $a=0{,}603$, $b=0{,}397$.
\end{enumerate}
D'où $\Delta = 0{,}5\cdot0{,}603 - 0{,}5\cdot0{,}397 = 0{,}103$ et
\[
  \alpha^\star(S) = \frac{0{,}103\,S}{0{,}5\times0{,}397} = 0{,}519\,S.
\]

\section{Validation empirique}

Balayage $S,\alpha\in\{0{,}1;0{,}3;0{,}35;0{,}5;0{,}7;0{,}9;1{,}0\}$, $M=200$,
$N=50$, $\mathtt{total}=1500$, 20 runs (\texttt{random\_stump\_perclause\_boost.cpp}).

\begin{table}[h]
\centering
\renewcommand{\arraystretch}{1.2}
\begin{tabular}{cccccccc}
\toprule
$S \backslash \alpha$ & 0,1 & 0,3 & 0,35 & 0,5 & 0,7 & 0,9 & 1,0 \\
\midrule
0,1 & 78,1 & \cellcolor{red!15}48,9 & \cellcolor{red!15}48,9 & \cellcolor{red!15}48,9 & \cellcolor{red!15}48,9 & \cellcolor{red!15}48,9 & \cellcolor{red!15}48,9 \\
0,3 & 86,1 & 74,2 & 64,3 & \cellcolor{red!15}48,9 & \cellcolor{red!15}48,9 & \cellcolor{red!15}48,9 & \cellcolor{red!15}48,9 \\
0,5 & 85,4 & 91,5 & 82,7 & 61,6 & \cellcolor{red!15}48,9 & \cellcolor{red!15}48,9 & \cellcolor{red!15}48,9 \\
0,7 & 84,1 & 95,8 & 89,0 & 77,0 & 51,4 & \cellcolor{red!15}48,9 & \cellcolor{red!15}48,9 \\
0,9 & 81,9 & 97,5 & 95,8 & 80,1 & 69,4 & \cellcolor{red!15}48,9 & \cellcolor{red!15}48,9 \\
1,0 & 81,1 & 98,6 & \textbf{99,1} & 83,4 & 75,8 & 56,5 & \cellcolor{red!15}48,9 \\
\bottomrule
\end{tabular}
\caption{Accuracy (\%, 20 runs). Fond rouge : $48{,}9\,\%$, variance nulle ---
effondrement total vers un classifieur trivial.}
\end{table}

\begin{table}[h]
\centering
\begin{tabular}{ccc}
\toprule
$S$ & $\alpha^\star$ prédit & Dernier $\alpha$ non-effondré (empirique) \\
\midrule
0,1 & 0,052 & 0,1 (78,1\,\%, déjà dégradé) \\
0,3 & 0,156 & 0,3 (74,2\,\%) \\
0,5 & 0,259 & 0,35 (82,7\,\%) \\
0,7 & 0,363 & 0,35 (89,0\,\%) \\
0,9 & 0,467 & 0,5 (80,1\,\%) \\
1,0 & 0,519 & 0,5 (83,4\,\%) \\
\bottomrule
\end{tabular}
\caption{Seuil prédit vs position réelle de l'effondrement. Le dernier $\alpha$
non-effondré est toujours juste sous $\alpha^\star(S)$, l'effondrement (48,9\,\%)
commence toujours juste au-dessus.}
\end{table}

\section{Conclusion}

Le meilleur point trouvé ($S{=}1{,}0,\alpha{=}0{,}35\to99{,}1\,\%$) est celui qui
maximise la marge sous $\alpha^\star(S)=0{,}519$ sans réduire $\alpha$ au point de
ralentir la correction du bruit ($\alpha=0{,}1$ plafonne à 78--86\,\% quel que
soit $S$, trop faible pour purger vite les mauvais états initiaux). La formule
$\alpha^\star(S)$ prédit correctement, sur les 42 cellules testées, où
l'effondrement total commence.

% ======================================================
\section{Étude de cas : Iris}
% ======================================================

\subsection{Différence de départ avec XOR}

Sur XOR, aucune feature seule ne prédit $y$ ($a\approx b$ pour toutes) : il faut
en combiner deux. Sur Iris, plusieurs features prédisent $y$ \textbf{seules},
avec un écart $a-b$ très supérieur à celui de XOR ($0{,}206$).

\subsection{Mesure de $c,a,b$ sur Iris (\texttt{BinaryIrisData.txt}, 16 features, 149 exemples)}

Pour chaque feature et chaque classe (one-vs-rest), on mesure $c=P(x{=}1)$,
$a=P(y{=}1\mid x{=}1)$, $b=P(y{=}1\mid x{=}0)$ et on calcule $\alpha^\star(S{=}1)$
avec la même formule que pour XOR :

\begin{table}[h]
\centering
\renewcommand{\arraystretch}{1.2}
\begin{tabular}{cccccc}
\toprule
Feature & Classe & $c$ & $a$ & $b$ & $\alpha^\star(S{=}1)$ \\
\midrule
8  & 1 & 0,631 & 0,468 & 0,091 & $\approx0$ (fragile) \\
11 & 2 & 0,624 & 0,484 & 0,089 & 0,60 (fragile) \\
2  & 2 & 0,282 & 0,738 & 0,178 & 4,05 \\
7  & 0 & 0,369 & 0,673 & 0,138 & 4,85 \\
11 & 1 & 0,624 & 0,516 & 0,018 & 10,00 \\
16 & 2 & 0,349 & 0,904 & 0,031 & 29,67 (quasi incassable) \\
\bottomrule
\end{tabular}
\caption{7 des 16 features sont constantes (toujours la même valeur, aucun
signal possible). Parmi les 9 restantes, la force du signal varie de très
fragile ($\alpha^\star\approx0$) à quasi incassable ($\alpha^\star\approx30$).}
\end{table}

\subsection{Validation empirique}

Balayage $S,\alpha\in\{0{,}05;0{,}1;0{,}3;0{,}5;0{,}7;1{,}0\}$, $M=1500$, $K=2$,
$N=100$, $\mathtt{total}=20\,000$, 20 runs
(\texttt{random\_stump\_multifeature\_multiclass.cpp}).

\begin{table}[h]
\centering
\renewcommand{\arraystretch}{1.2}
\begin{tabular}{ccccccc}
\toprule
$S \backslash \alpha$ & 0,05 & 0,1 & 0,3 & 0,5 & 0,7 & 1,0 \\
\midrule
0,05 & 95,5 & 92,0 & 89,8 & 88,0 & 88,0 & 88,2 \\
0,1  & 95,5 & 95,7 & 90,7 & 90,7 & 89,3 & 88,0 \\
0,3  & 95,5 & 95,5 & 95,3 & 90,8 & 92,0 & 91,8 \\
0,5  & 95,5 & 95,5 & 95,5 & 94,8 & 92,8 & 92,2 \\
0,7  & 95,5 & 95,5 & 95,5 & 95,5 & 95,2 & 92,8 \\
1,0  & 95,5 & 95,5 & 95,5 & 95,5 & 95,5 & 95,7 \\
\bottomrule
\end{tabular}
\caption{Accuracy (\%, 20 runs). Contrairement à XOR : jamais d'effondrement
total (le pire reste à 88\,\%, largement au-dessus du hasard à 33\,\%).}
\end{table}

\subsection{Conclusion}

Deux différences avec XOR expliquent le tableau :
\begin{itemize}
  \item La plupart des features/classes ont $\alpha^\star \gg 1$ : hors de la
  plage valide $(0,1]$, la destination ne bascule jamais --- d'où la zone plate
  à $95{,}5\,\%$ dès que $S$ est assez grand.
  \item La dégradation à petit $S$/grand $\alpha$ (95,5\,\%$\to$88\,\%) n'est
  donc pas un basculement de destination : c'est un effet de \emph{vitesse}
  (convergence pas assez rapide dans les 20\,000 tirages), comme la colonne
  $\alpha=0{,}1$ de XOR (Section~5). Les quelques features fragiles
  ($\alpha^\star\approx0{,}6$, dans la plage testée) existent, mais leur effet
  se dilue dans le vote de 1500 clauses tirées parmi 120 paires de features
  possibles.
\end{itemize}
XOR est proche du seuil critique (bascule possible et observée) ; Iris en est
loin (jamais de bascule, seulement une question de vitesse).

% ======================================================
\section{Étude de cas : Digits}
% ======================================================

\subsection{Mesure de $c,a,b$ sur Digits (\texttt{BinaryDigitsTrainingData192.txt},
192 features, 1437 exemples, 10 classes)}

Sur $192\times10=1920$ combinaisons feature/classe, $312$ ont un signal IDENTITY
exploitable ($|a-b|\ge0{,}02$). Contrairement à Iris, plusieurs sont
\textbf{très fragiles} --- $\alpha^\star(S{=}1)$ tombe sous $1$ pour un nombre
significatif de features, pas seulement 1 ou 2 cas isolés comme sur Iris :

\begin{table}[h]
\centering
\renewcommand{\arraystretch}{1.2}
\begin{tabular}{cccccc}
\toprule
Feature & Classe & $c$ & $a$ & $b$ & $\alpha^\star(S{=}1)$ \\
\midrule
55  & 0 & 0,511 & 0,125 & 0,075 & 0,113 (très fragile) \\
26  & 6 & 0,505 & 0,121 & 0,081 & 0,293 \\
47  & 9 & 0,522 & 0,143 & 0,048 & 0,333 \\
55  & 9 & 0,511 & 0,129 & 0,064 & 0,378 \\
18  & 9 & 0,502 & 0,121 & 0,074 & 0,547 \\
180 & 1 & 0,500 & 0,132 & 0,082 & 0,593 \\
\midrule
34  & 4 & 0,410 & 0,219 & 0,007 & 63,67 \\
34  & 0 & 0,410 & 0,241 & 0,004 & 132,67 \\
31  & 0 & 0,390 & 0,255 & 0,002 & 229,00 \\
\bottomrule
\end{tabular}
\caption{Échantillon des cas les plus fragiles (haut) et les plus robustes (bas)
parmi les 312 combinaisons exploitables. Contrairement à Iris (une seule feature
fragile trouvée, $\alpha^\star{=}0{,}60$), Digits en a plusieurs sous $1$.}
\end{table}

\textbf{Prédiction avant validation} (voir Section~7.3 pour le résultat réel,
qui infirme cette prédiction) : plus de features fragiles dans la plage
testable $\Rightarrow$ on s'attend à une dégradation plus marquée que sur Iris à
petit $S$/grand $\alpha$, potentiellement plus proche du comportement de XOR.

\subsection{Validation empirique}

Balayage $S,\alpha\in\{0{,}05;0{,}1;0{,}3;0{,}5;0{,}7;1{,}0\}$, $M=1000$, $K=2$,
$N=100$, $\mathtt{total}=20\,000$, 2 runs
(\texttt{random\_stump\_multifeature\_multiclass.cpp}).

\begin{table}[h]
\centering
\renewcommand{\arraystretch}{1.2}
\begin{tabular}{ccccccc}
\toprule
$S \backslash \alpha$ & 0,05 & 0,1 & 0,3 & 0,5 & 0,7 & 1,0 \\
\midrule
0,05 & 95,6 & 95,8 & 96,5 & 96,4 & 96,4 & \textbf{97,1} \\
0,1  & 95,6 & 95,6 & 95,8 & 96,5 & 96,9 & 96,5 \\
0,3  & 95,6 & 95,6 & 95,4 & 96,0 & 96,0 & 96,8 \\
0,5  & 95,6 & 95,6 & 95,4 & 95,6 & 95,0 & 96,3 \\
0,7  & 95,6 & 95,6 & 95,4 & 95,6 & 95,3 & 95,0 \\
1,0  & 95,6 & 95,6 & 95,4 & 95,6 & 95,6 & 94,2 \\
\bottomrule
\end{tabular}
\caption{Accuracy (\%, 2 runs). Aucun effondrement nulle part : tout reste entre
94,2\,\% et 97,1\,\%, plus stable qu'Iris (qui descendait jusqu'à 88\,\%).}
\end{table}

\subsection{Conclusion --- la prédiction était fausse, et pourquoi}

La prédiction de la section précédente ne se vérifie pas : Digits est en
réalité \textbf{plus stable} qu'Iris, pas moins, malgré des features
individuellement plus fragiles ($\alpha^\star=0{,}113$ contre $0{,}60$ pour le
pire cas d'Iris).

L'erreur venait d'un raisonnement incomplet : la fragilité d'une feature
\emph{seule} ne suffit pas à prédire l'effet sur le vote global --- il faut
aussi tenir compte du nombre de combinaisons possibles. Digits a $192$
features, donc $\binom{192}{2}\approx18\,300$ paires conditionnables, contre
$\binom{16}{2}=120$ pour Iris (environ $150\times$ plus). Même si quelques
features sont fragiles, la probabilité qu'une clause tire précisément une
combinaison qui isole l'une d'elles comme littéral résiduel est bien plus
faible sur Digits --- l'effet se dilue davantage dans le vote de milliers de
clauses, malgré une fragilité individuelle plus forte.

\begin{note}[Leçon méthodologique, corrigée après le test MNIST --- voir Section~8]
Le nombre de combinaisons possibles seul ne suffit pas à prédire la dilution :
MNIST (Section~8) a encore plus de combinaisons que Digits, et pourtant montre
\emph{beaucoup plus} de sensibilité à $S,\alpha$. Le vrai facteur est le
\textbf{ratio entre $M$ (nombre de clauses) et le nombre de combinaisons
possibles} $\binom{n}{K}$ --- pas la taille de l'espace seule. Voir Section~8.3.
\end{note}

% ======================================================
\section{Étude de cas : MNIST}
% ======================================================

\subsection{Mesure de $c,a,b$ sur MNIST (784 features, 60\,000 exemples, 10 classes)}

Sur $784\times10=7840$ combinaisons feature/classe, $1905$ ont un signal
IDENTITY exploitable. Le pire cas trouvé est encore plus fragile que sur
Digits : $\alpha^\star(S{=}1)=0{,}051$ (feature 434, classe 1), contre $0{,}113$
pour Digits.

\subsection{Validation empirique}

Balayage $S,\alpha\in\{0{,}05;0{,}1;0{,}3;0{,}5;0{,}7;1{,}0\}$, $M=500$
(réduit pour un balayage rapide), $K=2$, $N=100$, $\mathtt{total}=20\,000$,
1 run (\texttt{mnist\_onyxia}).

\begin{table}[h]
\centering
\renewcommand{\arraystretch}{1.2}
\begin{tabular}{ccccccc}
\toprule
$S \backslash \alpha$ & 0,05 & 0,1 & 0,3 & 0,5 & 0,7 & 1,0 \\
\midrule
0,05 & 85,9 & 80,5 & 58,4 & 52,7 & 51,0 & 39,6 \\
0,1  & 85,9 & 85,9 & 71,5 & 58,1 & 51,7 & 48,8 \\
0,3  & 85,9 & 85,9 & 85,9 & 77,4 & 71,7 & 63,7 \\
0,5  & 85,9 & 85,9 & 85,9 & 85,9 & 79,7 & 74,2 \\
0,7  & 85,9 & 85,9 & 85,9 & 85,9 & 85,9 & 78,8 \\
1,0  & 85,9 & 85,9 & 85,9 & 85,9 & 85,9 & 85,5 \\
\bottomrule
\end{tabular}
\caption{Accuracy (\%, 1 run). Forte dégradation (85,9\,\%$\to$39,6\,\%) quand
$\alpha$ dépasse $S$ --- contrairement à Digits, quasi plat.}
\end{table}

\subsection{Pourquoi MNIST est sensible alors que Digits ne l'est pas : le ratio M/combinaisons}

MNIST a \emph{plus} de combinaisons possibles que Digits
($\binom{784}{2}\approx307\,000$ contre $18\,300$), donc la dilution
« espace combinatoire seul » prédirait moins de sensibilité, pas plus. Le
facteur manquant : $M$ testé ici est \textbf{deux fois plus petit} pour MNIST
(500) que pour Digits (1000). Rapporté à la taille de l'espace :
\[
  \text{Digits : } \frac{1000}{18\,300}\approx5{,}5\,\%,
  \qquad
  \text{MNIST : } \frac{500}{307\,000}\approx0{,}16\,\% \quad (34\times\text{plus faible}).
\]
Avec un tirage aussi clairsemé, chaque clause pèse proportionnellement plus
dans le vote : les clauses qui tombent sur une combinaison fragile \emph{et}
un mauvais réglage ne sont plus diluées par la masse des autres. C'est cette
rareté relative de l'échantillonnage, pas la taille absolue de l'espace, qui
gouverne la sensibilité à $S,\alpha$.

\section{Synthèse}

\begin{center}
\renewcommand{\arraystretch}{1.3}
\begin{tabular}{lcccc}
\toprule
Dataset & $\binom{n}{K}$ & $M$ testé & $M/\binom{n}{K}$ & Sensibilité \\
\midrule
XOR    & $12$      & 200  & $17\,\%$    & Effondrement net (mur à $\alpha^\star$) \\
Iris   & $120$     & 1500 & $1250\,\%$  & Dégradation modérée (88--95,5\,\%) \\
Digits & $18\,300$ & 1000 & $5{,}5\,\%$ & Quasi aucune (94,2--97,1\,\%) \\
MNIST  & $307\,000$& 500  & $0{,}16\,\%$& Forte (39,6--85,9\,\%) \\
\bottomrule
\end{tabular}
\end{center}

La colonne qui explique vraiment la sensibilité est $M/\binom{n}{K}$ (le
« taux d'échantillonnage » de l'espace combinatoire), pas $\binom{n}{K}$ seul :
Iris (ratio le plus haut) est robuste, MNIST (ratio le plus bas, malgré le plus
grand espace) est le plus sensible. Un $M$ plus grand sur MNIST (non testé ici,
piste à suivre) devrait réduire cette sensibilité en rapprochant le ratio de
celui de Digits.

\end{document}
